\chapter{Estado de la Cuestión}
\label{cap:estadoDeLaCuestion}

% En el estado de la cuestión es donde aparecen gran parte de las referencias bibliográficas del trabajo. Una de las formas más cómodas de gestionar la bibliografía en {\LaTeX} es utilizando \textbf{bibtex}. Las entradas bibliográficas deben estar en un fichero con extensión \textit{.bib} (con esta plantilla se proporciona el fichero biblio.bib, donde están las entradas referenciadas más abajo). Cada entrada bibliográfica tiene una clave que permite referenciarla desde cualquier parte del texto con los siguiente comandos:

% \begin{itemize}
% \item Referencia bibliografica con cite: \cite{ldesc2e}
% \item Referencia bibliográfica con citep: \citep{notsoshort}
% \item Referencia bibliográfica con citet: \citet{latexAPrimer}
% \end{itemize}

% Es posible citar más de una fuente, como por ejemplo \citep{latexCompanion,LaTeXLamport,texKnuth}

% Después, latex se ocupa de rellenar la sección de bibliografía con las entradas \textbf{que hayan sido citadas} (es decir, no con todas las entradas que hay en el .bib, sino sólo con aquellas que se hayan citado en alguna parte del texto).

% Bibtex es un programa separado de latex, pdflatex o cualquier otra cosa que se use para compilar los .tex, de manera que para que se rellene correctamente la sección de bibliografía es necesario compilar primero el trabajo (a veces es necesario compilarlo dos veces), compilar después con bibtex, y volver a compilar otra vez el trabajo (de nuevo, puede ser necesario compilarlo dos veces). 

% \subsection{El giro computacional en el análisis de textos artísticos y creativos}
% (TODO: revisar que flipas)
% El procesamiento de lenguaje natural (PLN) ha evolucionado rápidamente en el transcurso de los últimos años. Aunque inicialmente [...], este desarrollo ha alcanzado la mejora en el ámbito de la generación de contenido creativo e incluso poético. 




% El desarrollo de una arquitectura capaz de automatizar el análisis, catalogación e interrogación de grandes corpus de música folclórica mediante Inteligencia Artificial requiere converger disciplinas históricamente distantes. Este capítulo examina el marco conceptual y los antecedentes científicos organizados en tres ejes fundamentales: el giro computacional en las humanidades y el procesamiento simbólico de partituras, la emergencia de las tecnologías basadas en Modelos de Lenguaje de Gran Escala (LLMs), y los paradigmas contrapuestos de recuperación de información simbólica y relacional.

% \section{El giro computacional en el análisis de textos artísticos y el procesamiento simbólico musical}
% \label{sec:giro_computacional}

% En las últimas décadas, las disciplinas humanísticas han experimentado lo que la literatura científica denomina el \textit{``giro computacional''} (\textit{computational turn}), un desplazamiento metodológico que transita desde el análisis microscópico de obras aisladas hacia la explotación macroscópica de grandes volúmenes de datos mediante técnicas de \textit{Distant Reading} \citep{moretti2013distant}. En el dominio de la musicología, este cambio de paradigma se ha consolidado bajo la disciplina de la musicología computacional, habilitando el estudio de patrones estilísticos, estructuras melódicas y evoluciones folclóricas sobre colecciones que exceden la capacidad de absorción del análisis manual analógico \citep{cook2004computational}.

% El éxito de este enfoque está íntimamente ligado a la estandarización de formatos de representación simbólica digital. Estructuras lógicas y lenguajes de marcado como \textit{MusicXML} y, más recientemente, el estándar de la \textit{Music Encoding Initiative} (MEI), han permitido codificar no solo las propiedades acústicas o de rendimiento de una obra (a la manera del formato MIDI estándar), sino su semántica textual y ontología musical compleja, integrando elementos líricos, métricos y estructurales dentro de esquemas arbóreos jerárquicos basados en XML \citep{roland2002music}.

% No obstante, la manipulación directa de estos ficheros impone una elevada complejidad algorítmica. Para superar esta barrera, la comunidad científica ha desarrollado \textit{frameworks} orientados a objetos específicos para la manipulación musicológica. Entre ellos, el ecosistema \textbf{music21} se ha consolidado como el estándar de facto en el ámbito de la investigación \citep{cuthbert2010music21}. Concebido como un conjunto de herramientas de código abierto basadas en Python, \textit{music21} abstrae la complejidad de los formatos de marcado subyacentes, permitiendo a los investigadores programar algoritmos analíticos para la búsqueda de patrones de intervalos, la extracción de líneas melódicas y el cómputo de jerarquías métricas. A pesar de su robustez, el uso de estos entornos sigue requiriendo competencias avanzadas de programación, lo que introduce una brecha técnica crítica para los investigadores de corte humanístico y musicológico tradicional.

% \section{Modelos de Lenguaje de Gran Escala (LLMs) y el dominio musical simbólico}
% \label{sec:llms_musicales}

% La llegada de la arquitectura \textit{Transformer} \citep{vaswani2017attention} revolucionó el procesamiento del lenguaje natural (NLP) y, de forma paulatina, ha comenzado a influir en las metodologías de análisis de datos estructurados y creativos. La capacidad intrínseca de los Modelos de Lenguaje de Gran Escala (LLMs) para modelar dependencias a largo plazo y capturar representaciones semánticas complejas ha impulsado el desarrollo de iniciativas recientes orientadas a la comprensión y generación de música simbólica.

% Trabajos pioneros en este campo exploraron el entrenamiento de modelos de lenguaje autorregresivos sobre representaciones textualizadas de partituras (como notación ABC o secuencias linearizadas de eventos MIDI). Sin embargo, la investigación contemporánea ha detectado limitaciones severas en estos enfoques debido a la naturaleza multidimensional y polifónica de la música, características que la sintaxis puramente secuencial del lenguaje de texto no puede codificar de manera directa \citep{ji2023comprehensive}.

% Recientemente, el paradigma ha evolucionado hacia sistemas multimodales capaces de alinear codificadores musicales especializados con modelos fundacionales de texto. Un ejemplo representativo es la emergencia de arquitecturas de instrucción como \textit{MIDI-LLaMA} o adaptaciones basadas en \textit{MusicBERT}, las cuales buscan que el modelo comprenda estructuras armónicas latentes mediante técnicas de alineamiento de características y ajuste fino conversacional \citep{chou2024midillama}. A pesar de estos avances de frontera, los modelos generales de última generación (como los pertenecientes a las familias \textit{Gemini} o \textit{Llama}) siguen manifestando deficiencias estructurales cuando operan de manera directa sobre partituras complejas codificadas en lenguaje natural: la densidad sintáctica de los formatos XML satura con rapidez las ventanas de contexto operativo y degrada los mecanismos de atención, provocando fallos sistemáticos en tareas que involucran razonamiento musical abstracto, cálculos matemáticos métricos o análisis de fenómenos expresivos sutiles como síncopas y hemiolias.

% \section{Recuperación de Información en Grandes Corpus: RAG Semántico frente a Enfoques Relacionales (Text-to-SQL)}
% \label{sec:rag_vs_sql}

% Ante la imposibilidad técnica de inyectar colecciones masivas de datos musicales en el contexto de un modelo de lenguaje, la ingeniería de la IA ha bifurcado sus esfuerzos de recuperación de información en dos aproximaciones distintas:

% \subsection{Generación Aumentada por Recuperación (RAG)}
% El enfoque RAG (\textit{Retrieval-Augmented Generation}) se basa en segmentar el corpus de documentos, proyectar dichos fragmentos en un espacio vectorial continuo mediante modelos de \textit{embeddings} y recuperar de forma semántica las secciones más relevantes de cara a fundamentar la respuesta del generador \citep{lewis2020retrieval}. Implementaciones avanzadas proporcionadas por la industria, tales como el motor nativo \textit{File Search} de Google Gemini, han facilitado el despliegue de estos sistemas sobre repositorios documentales extensos. No obstante, en contextos científicos de alta especificidad (como la musicología técnica), el RAG tradicional vectorizado o basado en texto presenta serios problemas de fiabilidad. Al carecer de capacidades intrínsecas de cómputo estructurado, los modelos de lenguaje tienden a la generación de fenómenos de alucinación, deduciendo métricas inexistentes o inventando referencias líricas o de autoría ausentes en el conjunto original de datos \citep{huang2023citation}.

% \subsection{Aproximaciones Relacionales y Text-to-SQL}
% Con el fin de limitar la variabilidad de los modelos conversacionales puros, la literatura científica ha propuesto el paradigma de traducción de lenguaje natural a lenguaje de consultas estructuradas (\textit{Text-to-SQL}). En este diseño de arquitectura, el modelo de lenguaje no actúa como el almacén o el recuperador indexado de los datos, sino como un agente traductor (un sintetizador de consultas) intercalado entre la interfaz de usuario y un motor de base de datos relacional determinista como \textit{SQLite} \citep{katsogiannis2023text}. Almacenar las variables musicológicas extraídas previamente mediante un procesamiento algorítmico estricto (por ejemplo, computadas previamente con \textit{music21}) garantiza que la fase de recuperación de información sea determinista, auditable y libre de alucinaciones factuales. Sin embargo, este enfoque desplaza la fuente de error hacia la capacidad de inferencia del LLM: si el modelo comete fallos sintácticos al construir la \textit{query} SQL, o si es incapaz de plasmar una condición musicológica abstracta dentro de las restricciones de un esquema de base de datos relacional, el sistema comete fallos críticos de recuperación.

% \subsection{Integración de Contexto y el Protocolo MCP}
% Como respuesta a la necesidad de estandarizar la conexión segura, bidireccional y eficiente entre agentes de IA y fuentes de datos relacionales locales (o APIs externas), la industria del software ha introducido frameworks de conectividad avanzados como el \textbf{Model Context Protocol (MCP)} impulsado a finales de 2024 \citep{anthropic2024mcp}. MCP define un protocolo abierto cliente-servidor que formaliza cómo un modelo de lenguaje puede descubrir y ejecutar herramientas y recursos específicos en su entorno operativo. En lugar de desarrollar conectores acoplados para cada base de datos, el protocolo MCP permite desacoplar la base de datos (por ejemplo, encapsulando las operaciones de extracción \textit{SQLite}) en un servidor independiente con el que el modelo interactúa dinámicamente según el contexto de la consulta planteada por el usuario. La investigación sobre la viabilidad de este protocolo combinado con modelos locales pequeños (SLMs, como \textit{Llama 3.1:8B}) para la ejecución de pipelines científicos específicos y la construcción de aplicaciones web guiadas representa uno de los campos de desarrollo más activos y con mayor potencial de transferencia tecnológica en la ingeniería de software actual.



% \section{SEGUNDA VERSIÓN}

El desarrollo de una arquitectura capaz de automatizar el análisis, catalogación e interrogación de grandes corpus de música folclórica mediante Inteligencia Artificial requiere la convergencia de disciplinas históricamente distantes. Este capítulo examina el marco conceptual, la evolución historiográfica y los antecedentes científicos organizados en cuatro ejes fundamentales: la evolución del análisis filológico y musical antes y después del paradigma computacional, la estandarización del procesamiento simbólico de partituras, la emergencia cronológica de las tecnologías de procesamiento del lenguaje natural (PLN) aplicadas al dominio musical, y los paradigmas metodológicos contrapuestos para la recuperación de información y la orquestación de contextos.

\section{Evolución del análisis de textos y música: del enfoque analógico al giro computacional}
\label{sec:evolucion_analisis}

\subsection{El paradigma tradicional y el problema humano en la catalogación}
% Históricamente, tanto la filología como la musicología se han apoyado en metodologías cualitativas orientadas al análisis en profundidad de casos concretos mediante la interpretación de textos aislados y el estudio detallado de partituras individuales. En el ámbito de la etnomusicología y el folklore, los esfuerzos pioneros de coleccionismo y catalogación del siglo XIX y principios del XX (tales como los cancioneros manuscritos y los primeros registros fonográficos) dependieron de metodologías estrictamente manuales de transcripción y etiquetado analógico \citep{bartok1951serbo}. 

% Este enfoque tradicional ha situado históricamente a las investigadoras ante un problema humano y operativo inevitable: la falta de escalabilidad. La catalogación sistemática de un corpus regional extenso exige una inversión de tiempo que a menudo excede la totalidad de una carrera académica. Tareas críticas como la indexación de variantes melódicas, la identificación de concordancias líricas intertextuales o el control de calidad estructural de miles de piezas tradicionales han dependido de la memoria y la capacidad de resistencia cognitiva de las musicólogas \citep{nettl2015study}. Como consecuencia, los proyectos con intereses historico bibliográficos se veían forzados a restringir su alcance a muestras minúsculas y sesgadas, limitando la validez estadística de las conclusiones sobre la evolución y mutación de las tradiciones orales.

Históricamente, tanto la filología como la musicología se han apoyado en metodologías cualitativas orientadas al análisis intensivo de documentos individuales, priorizando la interpretación contextualizada de textos, manuscritos y partituras concretas. Este paradigma interpretativo ha constituido durante siglos la base metodológica de las Humanidades, permitiendo desarrollar análisis de gran profundidad sobre obras particulares, autores específicos y tradiciones culturales delimitadas. En consecuencia, la producción de conocimiento dependía fundamentalmente de la lectura, comparación e interpretación manual realizada por especialistas, quienes actuaban simultáneamente como analistas, catalogadores y gestores de la información.

En el ámbito de la etnomusicología y los estudios folclóricos, esta aproximación se materializó en los grandes proyectos de recopilación y preservación desarrollados entre finales del siglo XIX y principios del XX. Investigadores como Béla Bartók y Zoltán Kodály realizaron extensas campañas de trabajo de campo destinadas a registrar, transcribir y clasificar repertorios musicales de tradición oral, sentando las bases de la musicología comparada moderna \citep{bartok1951serbo}. La posterior incorporación de tecnologías de grabación sonora permitió ampliar el volumen de material recopilado, pero los procesos de transcripción, catalogación y análisis continuaron dependiendo casi exclusivamente de procedimientos manuales \citep{seeger1958prescriptive}.

Este enfoque tradicional situó históricamente a las investigadoras ante una limitación estructural difícilmente evitable: la falta de escalabilidad. A medida que las colecciones documentales crecían, el volumen de información superaba progresivamente la capacidad humana de procesamiento. La catalogación sistemática de un corpus regional extenso exigía años o incluso décadas de trabajo especializado, mientras que tareas como la identificación de variantes melódicas, la detección de relaciones intertextuales, la comparación de motivos recurrentes o la verificación de metadatos dependían de la memoria, experiencia y resistencia cognitiva de los equipos de investigación \citep{nettl2015study}.

Las consecuencias metodológicas de esta limitación fueron significativas. En numerosos estudios, la imposibilidad práctica de analizar colecciones completas obligó a seleccionar subconjuntos reducidos de documentos o repertorios que actuaban como muestras representativas. Aunque esta estrategia permitía mantener la viabilidad de la investigación, introducía inevitablemente sesgos de selección y dificultaba la obtención de conclusiones generalizables sobre fenómenos culturales complejos. La reconstrucción de procesos de transmisión oral, evolución estilística o difusión geográfica de repertorios musicales quedaba condicionada por la capacidad material de los investigadores para examinar manualmente los documentos disponibles.

Esta problemática no fue exclusiva de la musicología. Durante gran parte del siglo XX, disciplinas como la historia, la lingüística, la filología o los estudios literarios se enfrentaron a desafíos similares derivados del crecimiento exponencial de los archivos y colecciones documentales. La progresiva digitalización de bibliotecas, hemerotecas y repositorios culturales incrementó todavía más la distancia entre el volumen de información disponible y la capacidad humana para analizarla exhaustivamente \citep{crane2006big, unesco2003charter}. Este fenómeno dio lugar a una creciente demanda de metodologías capaces de complementar el análisis cualitativo tradicional mediante procedimientos automáticos de exploración, clasificación y recuperación de información.

La tensión entre profundidad interpretativa y capacidad de análisis masivo constituye uno de los factores que impulsaron el surgimiento de las Humanidades Digitales y de los enfoques computacionales aplicados al patrimonio cultural. Como señalaría posteriormente Moretti mediante el concepto de \textit{distant reading} \citep{moretti2013distant}, el incremento constante del volumen documental obligaba a desarrollar nuevas estrategias metodológicas capaces de analizar colecciones de una escala imposible de abordar mediante la lectura o revisión individual de cada elemento. Este cambio de paradigma sentó las bases para la incorporación progresiva de técnicas computacionales en la investigación textual y musical, proceso que transformaría profundamente las posibilidades analíticas de las disciplinas humanísticas durante las décadas siguientes.

\subsection{El análisis de texto y lírica folclórica antes y después de los modelos de lenguaje}
% El análisis de las letras de las canciones folclóricas presenta una complejidad singular debido a su naturaleza fronteriza entre la literatura oral y el patrimonio lingüístico. Antes de la aparición del procesamiento del lenguaje natural moderno, el análisis computacional de estos textos se limitaba a aproximaciones léxicas cuantitativas básicas, tales como el cómputo de frecuencias de palabras (\textit{bag-of-words}), índices de concordancia simples o la búsqueda de palabras clave en contexto (KWIC). Estos métodos primitivos eran incapaces de capturar el sentido abstracto, las estructuras metafóricas, las variantes dialectales o las mutaciones semánticas características de la lírica tradicional \citep{jurafsky2000speech}.

% Con la aparición del PLN contemporáneo y, específicamente, la llegada de las representaciones vectoriales (\textit{word embeddings}) y las arquitecturas de aprendizaje profundo de carácter secuencial, el análisis textual experimentó una transformación radical. El procesamiento de la lírica folclórica transitó desde el recuento estático de vocablos hacia el modelado semántico del contexto. Las técnicas actuales permiten el reconocimiento de entidades nombradas, el análisis de sentimiento latente, la detección automática de tópicos literarios (\textit{Topic Modeling}) y el alineamiento computacional entre las variaciones textuales de una misma canción a lo largo de diferentes regiones geográficas \citep{blei2003latent}. No obstante, el verdadero desafío de la etnomusicología computacional contemporánea reside en la necesidad de tomar un enfoque híbrido: el PLN ya no puede operar de forma aislada, sino que debe alinearse con las características técnicas de la música simbólica para comprender la relación simbiótica entre el contenido lírico del texto y la estructura rítmica y melódica de la partitura.

El análisis de las letras de las canciones folclóricas presenta una complejidad singular debido a su naturaleza híbrida entre patrimonio lingüístico, literatura oral y expresión cultural. A diferencia de otros géneros textuales, las letras tradicionales suelen transmitirse mediante procesos de oralidad, reinterpretación y adaptación local, generando múltiples variantes de una misma composición a lo largo del tiempo y del espacio geográfico. Esta característica dificulta enormemente la aplicación de metodologías convencionales de análisis textual, ya que los cambios léxicos superficiales pueden ocultar significados, estructuras narrativas o funciones simbólicas prácticamente equivalentes.

Antes de la consolidación del Procesamiento del Lenguaje Natural (PLN) moderno, el análisis computacional de textos musicales se apoyaba principalmente en métodos estadísticos de representación léxica. Entre los enfoques más extendidos se encontraban los modelos de bolsa de palabras (\textit{bag-of-words}), el cálculo de frecuencias de términos, las listas de concordancias y los sistemas de búsqueda de palabras clave en contexto (\textit{Key Word in Context}, KWIC), ampliamente utilizados en lingüística de corpus y filología digital \citep{jurafsky2000speech}. Aunque estas técnicas permitían identificar patrones de vocabulario recurrentes o detectar términos característicos de determinadas colecciones, trataban las palabras como unidades independientes y carecían de mecanismos para representar relaciones semánticas, polisemia o dependencias sintácticas complejas.

Durante las primeras décadas del siglo XXI, la incorporación de métodos estadísticos más sofisticados permitió superar parcialmente estas limitaciones. Técnicas como la Indexación Semántica Latente (\textit{Latent Semantic Analysis}, LSA) \citep{deerwester1990indexing} y, posteriormente, los modelos probabilísticos de tópicos como \textit{Latent Dirichlet Allocation} (LDA) \citep{blei2003latent}, facilitaron la identificación automática de temas recurrentes dentro de grandes colecciones documentales. En el ámbito de los estudios folclóricos, estas herramientas comenzaron a utilizarse para explorar patrones narrativos, motivos temáticos y relaciones entre repertorios geográficamente distantes, permitiendo analizar colecciones de textos de una escala anteriormente inabordable.

El siguiente gran avance se produjo con la introducción de las representaciones distribuidas del lenguaje o \textit{word embeddings}. Modelos como \textit{Word2Vec} \citep{mikolov2013efficient} y \textit{GloVe} \citep{pennington2014glove} abandonaron la representación discreta de las palabras para proyectarlas en espacios vectoriales continuos donde la proximidad geométrica refleja similitud semántica. Esta innovación permitió capturar relaciones lingüísticas complejas, identificar variantes léxicas relacionadas y modelar contextos semánticos de forma mucho más rica que los enfoques basados exclusivamente en frecuencias. Paralelamente, las arquitecturas neuronales secuenciales, especialmente las redes neuronales recurrentes (RNN) y los modelos de memoria a largo plazo (LSTM), introdujeron la capacidad de procesar dependencias contextuales extensas dentro de secuencias textuales \citep{hochreiter1997long}.

La verdadera transformación, sin embargo, llegó con la aparición de los modelos basados en la arquitectura Transformer \citep{vaswani2017attention} y el desarrollo posterior de los Modelos de Lenguaje de Gran Escala (LLMs). A diferencia de las técnicas anteriores, estos sistemas no se limitan a representar palabras o frases aisladas, sino que construyen representaciones contextualizadas capaces de integrar simultáneamente información léxica, sintáctica, semántica y pragmática. Gracias a ello, resulta posible abordar tareas de elevada complejidad como el reconocimiento de entidades nombradas, la extracción automática de relaciones, la clasificación temática contextualizada, la detección de metáforas, el análisis de sentimiento y la generación de resúmenes interpretativos \citep{devlin2019bert}.

Para la investigación etnomusicológica, estas capacidades abren posibilidades inéditas. Los modelos actuales permiten identificar relaciones entre variantes textuales de una misma canción, detectar fórmulas poéticas recurrentes, reconstruir familias de transmisión oral y analizar patrones temáticos distribuidos a través de grandes corpus multirregionales. Además, la flexibilidad conversacional de los LLMs facilita la formulación de preguntas complejas en lenguaje natural, reduciendo significativamente las barreras técnicas de acceso a grandes colecciones documentales.

No obstante, la lírica constituye únicamente una de las dimensiones que conforman el fenómeno musical. Mientras que los avances recientes del PLN han mejorado notablemente la comprensión del contenido textual, la música incorpora simultáneamente estructuras melódicas, armónicas, rítmicas y formales que no pueden representarse adecuadamente mediante lenguaje natural. En consecuencia, uno de los principales retos de la etnomusicología computacional contemporánea consiste en desarrollar enfoques híbridos capaces de integrar la comprensión semántica de los textos con el análisis estructural de la música simbólica. Esta convergencia entre PLN y representación musical constituye actualmente una de las líneas de investigación más activas dentro de las Humanidades Digitales y la Inteligencia Artificial aplicada al patrimonio cultural \citep{zhao2023comprehensive, ji2023comprehensive}.

\subsection{El giro computacional (computational turn)}
% Hacia finales del siglo XX, las disciplinas humanísticas comenzaron a absorber formalmente las tecnologías informáticas, dando lugar a lo que la literatura científica acuñó como el \textit{``giro computacional''} (\textit{computational turn}) o la consolidación de las Humanidades Digitales \citep{moretti2013distant, berry2011computational}. Este fenómeno no supuso meramente la adopción de herramientas informáticas para optimizar tareas de almacenamiento o edición tipográfica, sino una reconfiguración profunda de las Humanidades Digitales.

% Dentro de este nuevo marco conceptual, la aportación teórica más disruptiva fue la formulación del concepto de \textit{Distant Reading} (Lectura Distante) impulsado por Franco Moretti \citep{moretti2013distant}. Este enfoque establece una confrontación directa con el \textit{Close Reading} (Lectura Cercana) tradicional, una metodología heredada de la filología del siglo XIX que exige una atención lineal y cualitativa sobre textos o piezas individuales. Ante la imposibilidad de abarcar sólo con la lectura humana el creciente volumen de documentos procedentes de los archivos digitalizados, el \textit{Distant Reading} propone alejarse deliberadamente del texto individual para adoptar una perspectiva de conjunto. Al agrupar y procesar simultáneamente miles de documentos, el análisis se desplaza hacia la detección de regularidades abstractas, patrones macro-estilísticos, transformaciones morfológicas y mutaciones culturales morfológicas que operan a nivel de sistema y que resultan intrínsecamente invisibles mediante el escrutinio aislado de una obra individual \citep{moretti2013distant}.

% En el dominio de la ciencia musical, este cambio paradigmático se materializó en la consolidación de la musicología computacional, un área interdisciplinar que aplica el formalismo matemático, el modelado estadístico y los principios de la recuperación de información sobre partituras digitalizadas \citep{cook2004computational}. Aplicado al estudio de la música folclórica y tradicional, el giro computacional redefine la investigación del patrimonio inmaterial. La partitura ya no se aborda únicamente como una instrucción de ejecución acústica, sino como una estructura de datos relacional y jerárquica susceptible de ser sometida a análisis estadísticos con múltiples variables. Esto permite a las investigadoras modelar matemáticamente la evolución de los repertorios, trazar mapas de difusión geográfica de variantes melódicas, calcular la entropía de los sistemas tonales regionales y aislar el ADN estilístico de una tradición oral a una escala y con un nivel de representatividad científica inédita en la musicología analógica.

A finales del siglo XX y comienzos del XXI, las disciplinas humanísticas comenzaron a incorporar de forma sistemática las tecnologías informáticas dentro de sus metodologías de investigación. Este proceso, conocido en la literatura como \textit{computational turn} o giro computacional, constituye uno de los cambios epistemológicos más significativos experimentados por las Humanidades desde la consolidación de los métodos histórico-críticos modernos \citep{berry2011computational}. Lejos de limitarse a la informatización de archivos o a la digitalización de documentos, este fenómeno implicó una transformación profunda de las preguntas de investigación, de las metodologías empleadas y de las escalas de análisis consideradas posibles dentro de disciplinas tradicionalmente interpretativas.

La consolidación de las Humanidades Digitales estuvo estrechamente ligada al crecimiento exponencial de repositorios digitales, bibliotecas electrónicas y colecciones patrimoniales accesibles en línea. Por primera vez en la historia, millones de documentos, manuscritos, partituras, grabaciones sonoras y materiales culturales pudieron almacenarse, consultarse y procesarse de forma automatizada \citep{schreibman2004companion, burdick2012digital}. Este incremento masivo de la disponibilidad de datos culturales generó nuevas oportunidades analíticas, pero también puso de manifiesto las limitaciones de los enfoques tradicionales basados exclusivamente en la interpretación manual de casos individuales.

Dentro de este nuevo marco conceptual, una de las contribuciones teóricas más influyentes fue la formulación del concepto de \textit{Distant Reading} propuesta por Franco Moretti \citep{moretti2013distant}. Este enfoque plantea una confrontación explícita con el paradigma del \textit{Close Reading}, heredado de la filología y la crítica literaria tradicionales, donde el análisis se concentra en la lectura detallada de textos individuales. Frente a la imposibilidad material de examinar manualmente colecciones compuestas por miles o millones de documentos, Moretti propone desplazar la atención desde la obra particular hacia el sistema cultural en su conjunto. El objetivo deja de ser interpretar exhaustivamente un único texto para identificar patrones agregados, regularidades históricas, transformaciones estilísticas y dinámicas de circulación cultural que únicamente emergen cuando el corpus se analiza a gran escala.

Esta perspectiva estimuló el desarrollo de nuevas metodologías cuantitativas para el estudio de fenómenos culturales. En el ámbito literario surgieron técnicas de estilometría, minería de textos, modelado temático y análisis de redes, mientras que disciplinas como la historia cultural comenzaron a incorporar métodos estadísticos y visualizaciones masivas de datos \citep{jockers2013macroanalysis, ramsay2011reading}. Paralelamente, el campo de las \textit{Cultural Analytics} impulsado por Lev Manovich defendió la aplicación de técnicas computacionales para estudiar patrones culturales a escalas anteriormente inalcanzables, combinando métodos procedentes de la estadística, la recuperación de información y la ciencia de datos con preguntas propias de las ciencias humanas \citep{manovich2020cultural}.

En el dominio musical, este cambio paradigmático favoreció la consolidación de la musicología computacional como área interdisciplinar dedicada al estudio cuantitativo de repertorios musicales mediante herramientas algorítmicas y modelos formales \citep{cook2004computational, conklin1995multiple}. La disponibilidad creciente de colecciones digitalizadas en formatos simbólicos como MIDI, MusicXML, Humdrum o MEI permitió comenzar a tratar las partituras no únicamente como objetos destinados a la interpretación sonora, sino también como conjuntos estructurados de datos susceptibles de análisis computacional \citep{roland2002music, selfridge1997humdrum}.

Aplicado específicamente al estudio de la música tradicional y el patrimonio inmaterial, el giro computacional transformó profundamente las posibilidades de investigación. La partitura pasó a concebirse como una representación formal de fenómenos musicales susceptibles de cuantificación, comparación y modelado estadístico. Gracias a ello, se hizo posible analizar simultáneamente miles de melodías, reconstruir procesos históricos de transmisión oral, detectar familias de variantes melódicas, estudiar la difusión geográfica de determinados rasgos estilísticos o caracterizar estadísticamente repertorios completos \citep{volk2012melodic}. Problemas que anteriormente requerían décadas de trabajo manual comenzaron a abordarse mediante algoritmos capaces de procesar grandes colecciones musicales en tiempos reducidos.

No obstante, el giro computacional no supuso la sustitución de las metodologías interpretativas tradicionales, sino la aparición de un paradigma híbrido en el que la capacidad explicativa de la investigación humanística se complementa con herramientas computacionales capaces de operar sobre volúmenes masivos de información. Esta convergencia entre interpretación y análisis algorítmico constituye actualmente uno de los pilares metodológicos de las Humanidades Digitales contemporáneas y sienta las bases para la incorporación posterior de técnicas avanzadas de inteligencia artificial y modelos de lenguaje en el estudio del patrimonio textual y musical.

\section{El procesamiento simbólico musical y los estándares de codificación}
\label{sec:procesamiento_simbolico}

El éxito del análisis musical automatizado está intrínsecamente sujeto a la existencia de formatos de representación simbólica digital que vayan más allá de la mera reproducción acústica de un archivo de audio o de la rigidez secuencial de un archivo MIDI estándar. Si bien el formato MIDI captura eventos físicos de ejecución (como el inicio y el final de una nota o su velocidad), carece por completo de semántica musicológica: no distingue entre una enarmonía (como Do sostenido y Re bemol), no codifica la estructura del compás ni permite la inserción jerárquica de la lírica textual o las anotaciones editoriales.

Para resolver este vacío, la comunidad científica desarrolló lenguajes de marcado estructurados fundamentados en XML. El primer estándar de adopción masiva en la industria fue \textit{MusicXML}, diseñado originalmente por Michael Good como un formato de intercambio centrado en la notación musical occidental y la representación gráfica de la partitura, permitiendo una migración fluida de datos entre editores comerciales \citep{good2001musicxml, good2011lessons}.

Sin embargo, en el ámbito de la investigación académica, la catalogación musicológica y las humanidades digitales, el formato predominante es el impulsado por la \textbf{Music Encoding Initiative (MEI)}. MEI es un estándar de código abierto, gestionado por la comunidad científica, que define un modelo conceptual y un esquema XML jerárquico diseñado específicamente para la edición y la preservación del patrimonio musical \citep{roland2002music}. A diferencia de MusicXML, MEI no se limita a describir cómo debe lucir gráficamente una partitura, sino que permite codificar la ontología y el contexto historicobibliográfico del documento. Un archivo MEI estructura de manera nativa tres niveles de información interconectados dentro de una sintaxis arbórea: la descripción bibliográfica detallada de la fuente (\textit{meiHead}), la representación de la estructura musical e interpretativa (\textit{music}), y la vinculación con variantes de texto, anotaciones críticas y reproducciones digitales. Esta riqueza estructural convierte a MEI en la herramienta idónea para la musicología computacional aplicada al folclore, permitiendo registrar las múltiples variantes y estados de conservación de una melodía tradicional.

Para manipular la densidad algorítmica de estos esquemas jerárquicos basados en XML, la comunidad de investigación ha adoptado de forma unánime el \textit{framework} orientado a objetos \textbf{music21} \citep{cuthbert2010music21}. Concebido como una librería de código abierto en Python, \textit{music21} abstrae la complejidad sintáctica de MusicXML y MEI, traduciendo los archivos en estructuras jerárquicas de objetos en memoria. Esto permite a los investigadores programar de forma determinista algoritmos analíticos para la extracción de descriptores complejos, tales como perfiles melódicos, densidades de eventos, contornos de intervalos, síncopas o hemiolias \citep{cuthbert2010music21}. Sin embargo, este entorno presenta una barrera de acceso técnica crítica: exige que el musicólogo posea competencias avanzadas en programación, ensanchando la brecha metodológica e impidiendo que el analista tradicional interactúe de forma dinámica y directa con el corpus sin intermediación técnica.

\section{Modelos de lenguaje de gran escala (LLMs) y el dominio musical simbólico}
\label{sec:llms_musicales}

\subsection{Contextualización histórica del PLN: de las Cadenas de Markov a la arquitectura Transformer}
La aplicación de técnicas lingüísticas al dominio musical no es un fenómeno reciente. La consideración de la música como un sistema formal con reglas sintácticas y gramaticales ha sido explorada desde mediados del siglo XX. Las primeras aproximaciones computacionales utilizaron modelos probabilísticos basados en Cadenas de Markov para modelar la probabilidad de transición entre notas y predecir estructuras melódicas o armónicas simples \citep{conklin1995multiple}. Durante las décadas siguientes, el PLN se apoyó en gramáticas formales libres de contexto y, con el auge del aprendizaje profundo a principios del siglo XXI, en redes neuronales recurrentes (RNN) y arquitecturas de memoria a largo plazo (\textit{Long Short-Term Memory}, LSTM), las cuales permitieron procesar secuencias musicales con una mejor retención de la dependencia temporal \citep{hochreiter1997long, mangal2019lstm}.

No obstante, el verdadero cambio de paradigma se produjo en 2017 con la introducción de la arquitectura \textit{Transformer} basada exclusivamente en mecanismos de auto-atención (\textit{self-attention}) \citep{vaswani2017attention}. Al eliminar la necesidad de un procesamiento secuencial recurrente y permitir la paralelización masiva, los Transformers revolucionaron el modelado del lenguaje humano y abrieron la puerta a los Modelos de Lenguaje de Gran Escala (LLMs). Al eliminar la necesidad de un procesamiento secuencial e iterativo (característico de las redes neuronales recurrentes (RNN) y las estructuras LSTM), los Transformers introdujeron la capacidad de ejecutar una paralelización masiva de los datos durante la etapa de entrenamiento. En el dominio de la musicología computacional, esta arquitectura revolucionó el modelado de estructuras simbólicas de larga duración. Mientras que los modelos previos sufrían de la degradación o el desvanecimiento del gradiente al intentar vincular eventos separados por centenares de compases, el mecanismo de auto-atención calcula la relevancia de todos los tokens de una obra de forma simultánea. Esto permite al sistema modelar la memoria musical a largo plazo, identificando con precisión correspondencias estructurales distantes como la reexposición de un tema en una forma sonata, la aparición de un estribillo en un romance tradicional o la variación de un \textit{leitmotiv}.

Asimismo, la flexibilidad de la arquitectura Transformer ha posibilitado una bifurcación metodológica en la investigación musicológica técnica en función del tipo de representación simbólica empleado. Por un lado, las arquitecturas basadas en codificadores puros (\textit{encoder-only}), inspiradas en modelos como BERT, han demostrado una enorme eficacia en tareas de comprensión e interrogación analítica profunda (\textit{downstream tasks}). Modelos especializados como \textit{MusicBERT} explotan esta capacidad para realizar anotaciones automáticas de armonía, clasificación de estilos folclóricos, segmentación de frases melódicas y control de calidad sobre errores de transcripción en los cancioneros digitales \citep{zeng2021musicbert}. El codificador analiza el contexto de la partitura, entendiendo qué notas anteceden y suceden a un evento musical concreto para deducir su función tonal o estructural. Por otro lado, las arquitecturas basadas en decodificadores puros (\textit{decoder-only}), de carácter autorregresivo, se han consolidado como el estándar para la generación y predicción secuencial de material sonoro, permitiendo modelar la improvisación y la evolución estilística de los repertorios mediante el cálculo probabilístico del próximo evento musical \citep{dhariwal2020jukebox}.

Finalmente, la asimilación del Transformer en la investigación etnomusicológica contemporánea ha obligado a reformular la gestión de la codificación posicional (\textit{positional encodings}). Dado que esta arquitectura carece de una noción de orden secuencial al procesar todos los tokens en paralelo, depende de vectores que inyectan la posición temporal de cada elemento. En música, donde el tiempo no es puramente lineal sino métrico y jerárquico (un pulso se organiza dentro de un compás, que a su vez se integra en una frase), las incrustaciones de posición convencionales del lenguaje humano resultan insuficientes. Los avances más recientes en este ámbito han desarrollado esquemas de atención relativa y codificaciones de tiempo compuesto para capturar simultáneamente la posición absoluta de la nota en el eje cronológico y su posición relativa dentro del marco métrico-rítmico de la pieza \citep{ji2023comprehensive}. Esta sofisticación permite a los Transformers identificar fenómenos expresivos abstractos de alta complejidad como las síncopas, las hemiolias o las polirritmias, permitiendo utilizar estos modelos como herramientas analíticas de primer orden para desentrañar las regularidades y singularidades del patrimonio musical inmaterial de tradición oral.

\subsection{Trabajos pioneros y la textualización de partituras}
Aprovechando las capacidades autorregresivas de los primeros LLMs, investigadores del ámbito de la música computacional exploraron el entrenamiento y \textit{fine-tuning} de estos modelos sobre representaciones textualizadas de partituras. Trabajos pioneros demostraron la viabilidad de codificar la música simbólica en formatos legibles por texto, tales como la notación tradicional ABC, cadenas linearizadas de eventos MIDI en texto (como los tokens \textit{Note-On} y \textit{Note-Off}), o lenguajes específicos de secuenciación como \textit{Humdrum kern} \citep{sturm2016music, oore2020this}. Estos enfoques permitieron abordar la composición y el análisis musical como tareas de predicción estadístico-lingüística de próximos tokens. 

Sin embargo, la literatura contemporánea ha identificado limitaciones severas en estas aproximaciones de textualización pura. A diferencia del texto escrito, que se organiza de manera unidireccional y lineal, la música es intrínsecamente multidimensional, jerárquica y polifónica; múltiples eventos sonoros independientes ocurren en un eje vertical simultáneo (armonía) mientras interactúan en un eje horizontal distributivo (melodía y ritmo) \citep{zhao2023comprehensive}. Forzar una estructura polifónica compleja dentro de una sintaxis textual estrictamente secuencial satura los mecanismos de atención de los modelos y degrada su capacidad de razonamiento conceptual.

Esta limitación se hizo especialmente evidente a medida que la comunidad investigadora comenzó a desplazar el foco desde la generación musical hacia tareas de análisis y recuperación de información. Mientras que la predicción autorregresiva de secuencias puede tolerar ciertas simplificaciones de la representación simbólica, las preguntas propias de la investigación musicológica suelen requerir razonamientos estructurales, comparaciones entre obras, identificación de patrones formales o extracción de métricas específicas. En estos escenarios, la pérdida de información derivada de la textualización secuencial se convierte en un obstáculo significativo para la fiabilidad de las respuestas generadas. Como consecuencia, comenzó a emerger una línea de investigación orientada no tanto a perfeccionar la textualización de las partituras, sino a explorar nuevas formas de integrar los modelos de lenguaje con representaciones musicales estructuradas y mecanismos externos de recuperación de conocimiento.

\subsection{Uso de LLMs en la investigación musicológica}
Las limitaciones observadas en los enfoques basados exclusivamente en la textualización de partituras han propiciado una redefinición progresiva del papel de los modelos de lenguaje dentro de la investigación musicológica. En lugar de considerar al LLM como un sistema encargado de interpretar directamente toda la información musical simbólica, las investigaciones más recientes tienden a emplearlo como una capa de interacción semántica situada entre el investigador y diversas fuentes especializadas de conocimiento musical. Los LLMs ofrecen el potencial único de unificar el análisis cualitativo y cuantitativo bajo un mismo entorno conversacional: entienden el lenguaje natural de las investigadoras humanas, comprenden el contexto cultural e histórico expresado en los metadatos y poseen una flexibilidad semántica capaz de responder a preguntas científicas imprevistas sin requerir que un programador humano rediseñe el software desde cero. El LLM se convierte en un mediador cognitivo, permitiendo al musicólogo interrogar colecciones complejas mediante lenguaje natural y abriendo las herramientas de la musicología computacional a usuarios sin conocimientos técnicos de código.

Para salvar la brecha polifónica, la investigación de frontera ha evolucionado recientemente hacia modelos especializados y arquitecturas de instrucción multimodal. Proyectos como \textit{MIDI-LLaMA} o adaptaciones basadas en \textit{MusicBERT} entrenan codificadores específicos que alinean las características latentes de la música simbólica con las incrustaciones de texto de los LLMs \citep{chou2024midillama, zeng2021musicbert}. No obstante, los modelos generales comerciales de frontera (como las familias \textit{Gemini} o \textit{Llama}) siguen manifestando deficiencias insalvables cuando operan directamente sobre la sintaxis nativa de archivos MusicXML o MEI complejos: la inyección masiva de estos strings XML devora las ventanas de contexto operativo (\textit{context window}) y genera alucinaciones sistemáticas ante cálculos de métricas abstractas o detección de síncopas, lo que exige explorar metodologías complementarias de recuperación y control de información.

\section{Paradigmas de recuperación de información y orquestación de contexto}
\label{sec:paradigmas_recuperacion}

Ante la inviabilidad de alojar colecciones completas de datos musicales dentro de la memoria de contexto de un LLM local o comercial, la ingeniería de la inteligencia artificial ha desarrollado múltiples técnicas de recuperación de información, las cuales presentan notables divergencias en términos de determinismo y fidelidad científica.

\subsection{El paradigma de Retrieval-Augmented Generation en modelos de lenguaje}
% El paradigma RAG tradicional particiona los documentos del corpus, genera incrustaciones vectoriales densas a través de un modelo de \textit{embeddings} y almacena dicha información en un índice vectorial continuo para recuperar las secciones semánticamente más cercanas a la consulta del usuario \citep{lewis2020retrieval}. Implementaciones comerciales contemporáneas de última generación, como el motor nativo \textit{File Search} integrado en la API de Google Gemini, automatizan este pipeline a gran escala en la nube. 

% Sin embargo, en disciplinas científicas rigurosas como la etnomusicología técnica, este enfoque basado en la aproximación semántica textual manifiesta severas deficiencias. Al carecer de capacidades aritméticas y lógicas intrínsecas, el modelo generativo es incapaz de calcular agregaciones estadísticas precisas a partir de fragmentos de texto flotantes (por ejemplo, computar el porcentaje exacto de canciones en tonalidad de Sol Mayor dentro del corpus), induciendo a alucinaciones factuales de alta gravedad científica donde el sistema inventa referencias líricas, atribuciones de autoría o métricas musicales inexistentes en la fuente original \citep{huang2023citation}.

La aparición del paradigma de \textit{Retrieval-Augmented Generation} (RAG) supuso un cambio significativo en el desarrollo de sistemas basados en modelos de lenguaje. Frente a los enfoques tradicionales, donde el conocimiento disponible para el modelo queda limitado a los parámetros aprendidos durante el entrenamiento, RAG introduce un mecanismo de recuperación dinámica que permite consultar fuentes documentales externas durante la inferencia \citep{lewis2020retrieval}. De este modo, el modelo deja de depender exclusivamente de su memoria paramétrica y puede fundamentar sus respuestas en información actualizada y específica del dominio de aplicación.

En su formulación clásica, el paradigma RAG parte los documentos del corpus en fragmentos de tamaño reducido (\textit{chunks}), genera representaciones vectoriales densas mediante modelos de \textit{embeddings} y almacena dichas representaciones en un índice vectorial. Cuando el usuario formula una consulta, esta se transforma en un vector semántico y se recuperan los fragmentos más próximos según una métrica de similitud en el espacio latente. Finalmente, los textos recuperados se incorporan al contexto del modelo generativo, que sintetiza una respuesta a partir de la información proporcionada \citep{lewis2020retrieval}. Implementaciones comerciales contemporáneas, como el servicio \textit{File Search} integrado en la API de Google Gemini, automatizan este proceso a gran escala y permiten construir sistemas conversacionales especializados sin necesidad de realizar \textit{fine-tuning} del modelo base.

No obstante, las hipótesis fundamentales que sustentan el éxito de RAG no siempre se mantienen en dominios altamente estructurados como la musicología computacional. El mecanismo de recuperación opera sobre similitud semántica entre fragmentos textuales, pero carece de una comprensión explícita de las relaciones lógicas, jerárquicas y cuantitativas presentes en las estructuras musicales. En consecuencia, aunque el sistema puede localizar descripciones relevantes de una obra o de un repertorio, encuentra dificultades cuando la consulta requiere efectuar cálculos exactos, agregaciones estadísticas o razonamientos sobre estructuras simbólicas complejas.

Esta limitación resulta especialmente problemática en contextos de investigación científica, donde la precisión de los resultados es un requisito fundamental. Preguntas aparentemente sencillas para una base de datos estructurada, como determinar el porcentaje exacto de canciones en tonalidad de Sol Mayor, identificar todas las obras que cumplen simultáneamente varios criterios musicales o calcular distribuciones estadísticas sobre el corpus, exigen operaciones deterministas que no pueden resolverse únicamente mediante recuperación semántica. En estos escenarios, el modelo tiende a inferir respuestas plausibles a partir de los fragmentos recuperados en lugar de ejecutar los cálculos necesarios, favoreciendo la aparición de errores factuales y alucinaciones documentadas en la literatura reciente \citep{huang2023citation}.

Por ello, diversos trabajos recientes han comenzado a explorar arquitecturas híbridas en las que los modelos de lenguaje se combinan con bases de datos estructuradas, motores de consulta y herramientas externas especializadas. Bajo este enfoque, el modelo conserva su capacidad para comprender lenguaje natural y generar explicaciones interpretables, mientras que las operaciones analíticas críticas son delegadas a sistemas deterministas capaces de garantizar la exactitud de los resultados.

\subsection{Aproximaciones relacionales deterministas (Text-to-SQL)}
% Para restringir de forma absoluta el espacio de alucinación del modelo, la literatura científica ha consolidado el paradigma de traducción de lenguaje natural a lenguaje de consultas estructuradas (\textit{Text-to-SQL}). En este diseño arquitectónico, los datos y las variables del corpus no se vectorizan a ciegas, sino que se extraen previamente de manera algorítmica determinista (empleando, por ejemplo, los pipelines de \textit{music21}) y se almacenan de forma estructurada en un sistema de gestión de bases de datos relacionales local y embebido como \textit{SQLite} \citep{katsogiannis2023text}. 

% Bajo este enfoque, el LLM deja de operar como un almacén de memoria indexado o probabilístico y pasa a actuar estrictamente como un agente traductor sintáctico. El modelo recibe la intención en lenguaje natural de la investigadora musicóloga y la traduce a una \textit{query} SQL unívoca que se ejecuta sobre el motor determinista local, garantizando que los datos devueltos sean matemáticamente exactos y auditables \citep{katsogiannis2023text}. El riesgo metodológico de esta solución se desplaza hacia la capacidad de inferencia del LLM: si el modelo genera un error de sintaxis SQL o malinterpreta las restricciones del esquema relacional, el sistema devuelve un error crítico de ejecución o una consulta vacía.

Como respuesta a las limitaciones observadas en los enfoques basados en recuperación semántica, la literatura reciente ha consolidado el paradigma de traducción de lenguaje natural a lenguaje de consultas estructuradas (\textit{Text-to-SQL}) como una alternativa orientada a la precisión, la verificabilidad y la ejecución determinista. Este enfoque se sitúa en continuidad con los sistemas de respuesta sobre bases de datos, cuya evolución ha sido ampliamente estudiada en la literatura reciente sobre comprensión de lenguaje natural aplicada a datos estructurados \citep{yu2018spider, zhong2017seq2sql}.

A diferencia de los sistemas RAG, donde la información se representa en espacios vectoriales y se recupera mediante similitud semántica \citep{lewis2020retrieval, gao2023ragsurvey}, el paradigma Text-to-SQL parte de una representación explícita y normalizada de los datos en forma de esquema relacional. Esta transición permite sustituir la aproximación probabilística del conocimiento por una estructura lógica formal, donde las operaciones sobre los datos son ejecutadas por un motor de bases de datos con semántica determinista.

En este contexto, los datos del corpus no son embebidos en un espacio latente, sino que son previamente extraídos, transformados y estructurados de forma determinista mediante pipelines específicos (por ejemplo, utilizando herramientas como \textit{music21} para la extracción de características musicales simbólicas). Posteriormente, dicha información se almacena en un sistema de gestión de bases de datos relacional embebido y local, como \textit{SQLite}, lo que permite garantizar propiedades de consistencia, integridad y trazabilidad sobre el conjunto de datos \citep{sqlite, cuthbert2010music21}.

Bajo este paradigma, el papel del modelo de lenguaje se redefine de manera sustancial. El LLM deja de actuar como un sistema de recuperación o síntesis basado en conocimiento implícito, y pasa a desempeñar el rol de interfaz semántica entre la usuaria y el sistema de consulta estructurado. Su función consiste en interpretar la intención expresada en lenguaje natural y traducirla a una consulta SQL formal, en línea con los enfoques de \textit{semantic parsing} aplicados a bases de datos \citep{lehman2022building, zhong2017seq2sql}. Posteriormente, dicha consulta es ejecutada de manera determinista por el motor de base de datos, lo que permite garantizar que los resultados obtenidos sean exactos, reproducibles y auditables.

Este desacoplamiento entre interpretación lingüística y ejecución lógica traslada el problema desde el dominio de la veracidad factual al dominio de la corrección sintáctica y semántica de la consulta generada. En otras palabras, el sistema no corre el riesgo de “inventar” datos que no existen en el corpus, como ocurre en los fenómenos de alucinación documentados en modelos generativos \citep{huang2023citation}, pero sí puede fallar en la fase de traducción si el modelo produce una consulta SQL incorrecta, incompleta o no alineada con el esquema relacional subyacente. Este tipo de errores ha sido ampliamente estudiado en la literatura sobre evaluación de sistemas Text-to-SQL, donde se distinguen métricas de exactitud de ejecución y de coincidencia sintáctica \citep{yu2018spider}.

Desde un punto de vista metodológico, esta propiedad resulta especialmente relevante en contextos de investigación científica y humanística computacional, donde la exactitud de los resultados y la capacidad de replicación son requisitos fundamentales. Asimismo, la estructura relacional permite realizar operaciones analíticas complejas (como agregaciones, filtrados múltiples o cálculos estadísticos sobre el corpus) mediante lenguajes formales con garantías matemáticas, en contraste con las limitaciones de consistencia observadas en enfoques puramente generativos o basados en recuperación semántica \citep{gao2023ragsurvey}.

En este sentido, el paradigma Text-to-SQL se posiciona como un complemento natural a los sistemas de recuperación aumentada, desplazando el énfasis desde la similitud semántica hacia la ejecución determinista de consultas estructuradas. Esta transición ha sido explorada en trabajos recientes sobre modelos de lenguaje con uso de herramientas (\textit{tool-augmented LLMs}) y sistemas neuro-simbólicos, donde el modelo actúa como un intérprete de intención mientras delega el razonamiento computacional en motores especializados \citep{yao2023react, schick2023toolformer, karpas2022mrkl}. Esta separación funcional constituye la base de las arquitecturas híbridas contemporáneas aplicadas a la integración de modelos de lenguaje con sistemas de bases de datos.

\subsection{Evolución reciente: el estándar Model Context Protocol (MCP)}
% Como respuesta directa a la necesidad de estandarizar la conexión bidireccional, segura y desacoplada entre los modelos de lenguaje y las infraestructuras locales de datos o APIs externas, la industria del software introdujo a finales de 2024 el denominado \textbf{Model Context Protocol (MCP)} \citep{anthropic2024mcp}. MCP formaliza un protocolo abierto cliente-servidor que opera de manera análoga al *Language Server Protocol* (LSP) empleado en entornos de desarrollo modernos. 

% En lugar de construir conectores monolíticos fuertemente acoplados para cada base de datos o API específica, el protocolo MCP permite empaquetar de forma aislada las herramientas y las capacidades de acceso a recursos (como las funciones de consulta e interrogación de un archivo binario \textit{SQLite}) dentro de un servidor MCP autónomo \citep{anthropic2024mcp}. El cliente de la aplicación web expone estas herramientas al contexto operativo del LLM de manera dinámica. El modelo de lenguaje descubre en tiempo real qué herramientas tiene a su disposición y las invoca según las necesidades de la consulta planteada por la usuaria. La investigación sobre la viabilidad, latencia y robustez de este protocolo integrado con Modelos de Lenguaje Pequeños (SLMs, como la familia \textit{Llama 3.1:8B}) operados bajo servidores locales eficientes como \textit{Ollama}, constituye uno de los campos de mayor vanguardia y con mayor potencial de transferencia tecnológica en el ámbito contemporáneo de la ingeniería de software y las Humanidades Digitales aplicadas.

La creciente adopción de modelos de lenguaje en entornos profesionales ha puesto de manifiesto una limitación recurrente de las arquitecturas basadas en herramientas externas: la ausencia de mecanismos estandarizados para conectar modelos, bases de datos, sistemas de archivos, APIs y servicios especializados. Históricamente, cada aplicación debía implementar integraciones específicas para cada recurso externo, generando ecosistemas fuertemente acoplados, difíciles de mantener y con escasa interoperabilidad. Como respuesta a esta problemática, a finales de 2024 se presentó el \textit{Model Context Protocol} (MCP), un protocolo abierto diseñado para estandarizar la comunicación entre modelos de lenguaje y herramientas externas mediante una arquitectura cliente-servidor desacoplada \citep{anthropic2024mcp}.

Conceptualmente, MCP puede entenderse como una extensión de las ideas introducidas previamente por los sistemas de uso de herramientas (\textit{tool-augmented language models}), donde el modelo deja de depender exclusivamente de su conocimiento paramétrico y adquiere la capacidad de interactuar con recursos externos para obtener información o ejecutar acciones \citep{schick2023toolformer, yao2023react}. Asimismo, el protocolo se alinea con la filosofía de arquitecturas neuro-simbólicas modulares como MRKL, en las que el modelo de lenguaje actúa como un orquestador que delega tareas especializadas en componentes externos más adecuados para resolverlas \citep{karpas2022mrkl}.

MCP formaliza esta interacción mediante un protocolo estandarizado inspirado parcialmente en la filosofía del \textit{Language Server Protocol} (LSP), ampliamente utilizado en entornos modernos de desarrollo software para desacoplar editores de código y servidores de análisis semántico. Bajo este enfoque, las capacidades disponibles para el modelo se encapsulan en servidores MCP independientes, cada uno de los cuales expone herramientas, recursos o acciones específicas a través de una interfaz homogénea. El cliente encargado de interactuar con el modelo descubre dinámicamente estas capacidades y las incorpora al contexto operativo durante la inferencia \citep{anthropic2024mcp}.

Esta aproximación ofrece ventajas significativas frente a las integraciones tradicionales desarrolladas. En lugar de construir conectores específicos para cada base de datos, API o repositorio documental, los desarrolladores pueden encapsular la funcionalidad en servidores MCP reutilizables e independientes de la aplicación final. Como consecuencia, la lógica de acceso a datos queda separada de la lógica conversacional, favoreciendo la mantenibilidad, escalabilidad y reutilización de los componentes software.

Desde la perspectiva de la interacción con modelos de lenguaje, el protocolo introduce un mecanismo de descubrimiento dinámico de herramientas. Durante la ejecución, el modelo recibe una descripción estructurada de las capacidades disponibles y puede decidir qué recurso utilizar en función de la consulta formulada por la usuaria. Este comportamiento se aproxima a los paradigmas contemporáneos de \textit{agentic AI}, donde el modelo no solo genera texto, sino que planifica secuencias de acciones, selecciona herramientas apropiadas y combina múltiples fuentes de información para resolver tareas complejas \citep{yao2023react, wang2024surveyagents}.

En el contexto específico de las Humanidades Digitales y la musicología computacional, MCP resulta especialmente útil porque permite integrar modelos de lenguaje con infraestructuras de datos locales sin exponer directamente la complejidad técnica subyacente. Recursos como bases de datos relacionales SQLite, colecciones documentales, repositorios MusicXML o herramientas de análisis musical pueden presentarse al modelo como capacidades abstractas accesibles mediante una interfaz uniforme. De este modo, el LLM puede interpretar consultas formuladas en lenguaje natural y delegar las operaciones analíticas en herramientas especializadas capaces de proporcionar resultados deterministas y verificables.

Aunque la investigación académica sobre MCP todavía se encuentra en una fase incipiente debido a su reciente aparición, el protocolo representa una evolución natural de las líneas de investigación sobre sistemas híbridos, recuperación aumentada, agentes basados en herramientas y arquitecturas neuro-simbólicas. Su adopción por parte de múltiples proveedores de modelos y plataformas de desarrollo sugiere que podría convertirse en un estándar de interoperabilidad para la próxima generación de aplicaciones basadas en inteligencia artificial, donde los modelos de lenguaje actuarán como interfaces universales de acceso a ecosistemas heterogéneos de conocimiento y computación \citep{anthropic2024mcp}.

\subsection{Técnicas más allá del RAG tradicional}
% Aunque el debate técnico se centra con frecuencia en la dualidad entre bases de datos relacionales y Generación Aumentada por Recuperación (RAG), la literatura en recuperación de información (\textit{Information Retrieval}, IR) abarca un abanico mucho más amplio de metodologías avanzadas de gestión de conocimiento:

% \begin{itemize}
%     \item \textbf{Recuperación Léxica y Keyword Matching (BM25):} Basada en representaciones probabilísticas de la frecuencia de términos e inversas de documentos. Es altamente eficiente para localizar metadatos textuales o líricas exactas, pero carece por completo de comprensión semántica o musical conceptual \citep{robertson2009probabilistic}.
    
%     \item \textbf{Gráficos de Conocimiento y Graph-RAG:} Consiste en estructurar el corpus en un gráfico de ontologías interconectadas (entidades y relaciones, como \textit{Obra $\rightarrow$ Pertenece a $\rightarrow$ Región $\rightarrow$ Comparte Estilo con $\rightarrow$ Otra Obra}). El LLM interroga el gráfico mediante lenguajes de consulta específicos como GraphQL o Cypher. Ofrece una excelente contextualización relacional e histórico-bibliográfica, pero su construcción e inferencia algorítmica imponen un coste computacional masivo \citep{edge2024from}.
    
%     \item \textbf{Modelos de Recuperación Híbridos y Re-ranking:} Combinan la búsqueda vectorial densa con la búsqueda léxica dispersa, ordenando los resultados mediante un modelo secundario de puntuación (\textit{Cross-Encoder Cross-Attention}). Optimiza la relevancia semántica textual, pero hereda la incapacidad del modelo para realizar operaciones matemáticas rigurosas sobre datos estructurados \citep{thakur2021beir}.
% \end{itemize}

Aunque gran parte del debate contemporáneo sobre sistemas basados en modelos de lenguaje se ha articulado en torno a la dicotomía entre recuperación aumentada mediante generación (RAG) y consulta sobre bases de datos estructuradas, la literatura en Recuperación de Información (\textit{Information Retrieval}, IR) ha desarrollado un conjunto mucho más amplio de metodologías para organizar, recuperar y explotar conocimiento especializado \citep{manning2008introduction}. Estas aproximaciones persiguen objetivos distintos (maximizar la relevancia semántica, mejorar la explicabilidad, capturar relaciones complejas o incrementar la precisión factual) y, en muchos casos, constituyen estrategias complementarias más que alternativas excluyentes.

\begin{itemize}
    \item \textbf{Recuperación léxica y búsqueda basada en palabras clave (BM25).} Los sistemas clásicos de recuperación documental continúan desempeñando un papel fundamental en numerosos motores de búsqueda. Modelos como BM25 estiman la relevancia de un documento a partir de la frecuencia de aparición de los términos consultados y de su distribución estadística dentro de la colección documental \citep{robertson2009probabilistic}. Su principal ventaja radica en su elevada eficiencia computacional y en su capacidad para localizar coincidencias exactas de términos, títulos, autores o fragmentos líricos. Sin embargo, estos sistemas carecen de mecanismos explícitos para capturar relaciones semánticas profundas, por lo que presentan dificultades cuando diferentes expresiones lingüísticas describen un mismo concepto musical.
    \item \textbf{Recuperación densa mediante embeddings neuronales.} El desarrollo de modelos de representación distribuida permitió superar parcialmente las limitaciones de los enfoques léxicos tradicionales. Técnicas como Dense Passage Retrieval (DPR) proyectan consultas y documentos en espacios vectoriales continuos donde la proximidad geométrica refleja similitud semántica \citep{karpukhin2020dense}. Este enfoque constituye la base de la mayoría de sistemas RAG contemporáneos \citep{lewis2020retrieval}. No obstante, su eficacia depende críticamente de la calidad de los embeddings y continúa enfrentando dificultades cuando la consulta requiere razonamientos lógicos, operaciones matemáticas o interpretaciones basadas en estructuras jerárquicas complejas.
    \item \textbf{Modelos híbridos y técnicas de re-ranking.} Con el objetivo de combinar las ventajas de la recuperación léxica y vectorial, numerosos sistemas contemporáneos emplean arquitecturas híbridas que fusionan ambos tipos de búsqueda. Los documentos candidatos recuperados son posteriormente reordenados mediante modelos neuronales de mayor capacidad, como los \textit{cross-encoders} basados en Transformer, capaces de evaluar con mayor precisión la relevancia contextual de cada resultado \citep{nogueira2019passage, thakur2021beir}. Estos enfoques han demostrado mejoras significativas en tareas de búsqueda abierta y recuperación documental, aunque siguen dependiendo de representaciones textuales y no resuelven de forma inherente los problemas asociados al razonamiento estructurado.
    \item \textbf{Gráficos de conocimiento y Graph-RAG.} Una línea de investigación particularmente activa consiste en representar la información mediante grafos semánticos donde las entidades y sus relaciones se almacenan explícitamente. Este enfoque permite modelar conexiones complejas entre obras musicales, autores, regiones geográficas, estilos o periodos históricos, facilitando consultas relacionales difíciles de expresar mediante recuperación vectorial convencional \citep{hogan2021knowledge}. La reciente aparición de arquitecturas conocidas como \textit{Graph-RAG} combina estos grafos con modelos de lenguaje para proporcionar mecanismos de recuperación más explicables y contextualizados \citep{edge2024from}. Sin embargo, la construcción, mantenimiento y actualización de grandes grafos de conocimiento requiere procesos de ingeniería considerablemente más costosos que los sistemas basados en índices vectoriales o bases de datos relacionales.
    \item \textbf{Sistemas multiagente y recuperación basada en herramientas.} Las investigaciones más recientes han comenzado a explorar arquitecturas donde múltiples agentes especializados colaboran para resolver consultas complejas. En estos sistemas, los modelos de lenguaje pueden seleccionar dinámicamente diferentes herramientas de búsqueda, bases de datos o recursos externos según las características de la tarea planteada \citep{yao2023react, wang2024surveyagents}. Este paradigma se encuentra estrechamente relacionado con los desarrollos recientes en protocolos de interoperabilidad como MCP y representa una de las direcciones más prometedoras para la integración de recuperación de información y razonamiento asistido por herramientas.
\end{itemize}

La coexistencia de estas metodologías pone de manifiesto que no existe una estrategia universalmente óptima para todos los escenarios de recuperación de información. En dominios caracterizados por datos textuales poco estructurados, los enfoques basados en recuperación semántica suelen ofrecer excelentes resultados. Sin embargo, cuando las consultas requieren precisión cuantitativa, trazabilidad metodológica y operaciones analíticas sobre atributos explícitos (como ocurre frecuentemente en la investigación musicológica computacional), las aproximaciones basadas en estructuras relacionales continúan ofreciendo ventajas significativas en términos de exactitud, reproducibilidad y verificabilidad de los resultados obtenidos.